% natural_join.tex
\begin{figure}[H]
	\centering
	\begin{tikzpicture}[
		table/.style={
			matrix of nodes,
			nodes={draw, minimum width=1.5cm, minimum height=0.6cm, align=center},
			column sep=-\pgflinewidth,
			row sep=-\pgflinewidth,
			nodes in empty cells
		},
		header/.style={fill=gray!20},
		>=Stealth
		]
		
		% Table R
		\matrix[table, label=above:{Relation $R$}] (R) {
			|[header]| A & |[header]| B \\
			1 & X \\
			2 & Y \\
			2 & U \\
			4 & V \\
		};
		
		% Table S
		\matrix[table, right=2.5cm of R, label=above:{Relation $S$}] (S) {
			|[header]| A & |[header]| C \\
			1 & Z \\
			2 & W \\
			3 & T \\
		};
		
		% Join Symbol
		\path (R.east) -- (S.west) node[midway] (join) {\Large $\bowtie$};
		
		% Result Table
		\matrix[table, below=1.5cm of join, label=above:{Result: $R \bowtie S$}] (Result) {
			|[header]| A & |[header]| B & |[header]| C \\
			1 & X & Z \\
			2 & Y & W \\
			2 & U & W \\
		};
		
		% Arrows tracking the match
		\draw[->, thick, dashed] (R.south) -- (Result.north west);
		\draw[->, thick, dashed] (S.south) -- (Result.north east);
		
		% Formula description
		\node[below=0.2cm of Result] {\small $\pi_{R \cup S}(\sigma_{R.A = S.A}(R \times S))$};
		
	\end{tikzpicture}
	\caption{Pictorial representation of Relational Algebra Natural Join}
	\label{fig:natural_join}
\end{figure}